\section{A strict feedback advantage over the full rarity class}
\label{sec:feedback-full-class}
\noindent\textbf{PROVED-ON-PAPER with exact rational verification.}
This strengthens the preceding two-profile example to the entire
global-rarity class for one finite-support magnitude law. It retains the
deterministic total cap of two flips; no expected-budget comparison or
asymptotic rate separation is asserted. The earlier certificate is preserved.

Let $\rho=1/4$ and
$\mu=\tfrac12\delta_{1/4}+\tfrac12\delta_{1/8}$.
Write $q_1=q(1/4)$, $q_2=q(1/8)$. The full admissible class is exactly
$q_1,q_2\ge0$, $q_1+q_2\le1/2$; the individual upper bounds by one
are automatic. Put
\[
 \theta=\tfrac12(q_1+q_2),\qquad
 A=\mathbb E[Dq(D)]=q_1/8+q_2/16.
\]
The feasible $(A,\theta)$ region is the triangle with vertices
\[
 (0,0),\quad(1/32,1/4),\quad(1/16,1/4).
\]
Equivalently $0\le A\le1/16$ and
$4A\le\theta\le\min\{8A,1/4\}$.

The noninteractive comparison class is exactly the one defined in the
preceding theorem: preselected independent local acquisition policies,
arbitrary shared and local seeds, and arbitrary nonlinear terminal
postprocessing, with a deterministic total cap of two flips and no
cross-coin acquisition feedback or outcome-dependent cancellation.

\begin{theorem}\label{feedback:full-class}
Over the entire class $\mathcal Q_{1/4}(\mu)$, the minimax MSE of the
noninteractive comparison class is strictly greater than $15475/10^6$.
An explicit two-flip feedback estimator has worst-case MSE strictly less
than $15443/10^6$. Its minimax squared-error advantage therefore exceeds
$32/10^6$.
\end{theorem}

\begin{proof}
One fresh bit has head probability $p=5/16+2A$.
In the word order $(TT,TH,HT,HH)$, the same-coin probabilities are
\[
 P_O(A)=(61/128-2A,\ 27/128,\ 27/128,\ 13/128+2A),
\]
whereas the two-fresh probabilities are
$P_N(A)=((1-p)^2,(1-p)p,p(1-p),p^2)$.
These formulas integrate out the hidden retained magnitude and sign;
the algorithm observes neither. In particular the two-flip laws depend
on $A$, while the target $\theta$ still varies over its feasible interval.

The feedback policy takes a fresh coin after the first tail and revisits
after the first head. Its four probabilities are
\[
 P_A(A)=((11/16-2A)^2,\ (11/16-2A)(5/16+2A),\
                       27/128,\ 13/128+2A).
\]
It takes exactly two charged flips. Define the fixed rational constant
$\pi=1074499095/2147483648$, and set
$a_x=P_A(0)_x$, $b_x=P_A(1/32)_x$ for its four outcomes. Its output table is
\[
 y_x=\frac14\frac{\pi b_x}{(1-\pi)a_x+\pi b_x}.
\]
In the displayed word order this is exactly
\[
 (y_{TT},y_{TH},y_{HT},y_{HH})=
 \left(\frac{26862477375}{237281040413},
       \frac{3223497285}{24696819223},
       \frac{1074499095}{8589934592},
       \frac{22564480995}{146053120736}\right).
\]
Although this table was constructed from the null and the admissible hard
alternative $(q_1,q_2)=(0,1/2)$, the same fixed table is used for every
$q\in\mathcal Q_{1/4}(\mu)$. It requires no unknown-parameter calibration.

Its risk is the polynomial
$R(A,\theta)=\sum_x P_A(A)_x(y_x-\theta)^2$.
For fixed $A$, its second derivative in $\theta$ is exactly two, so the
maximum over the vertical feasible interval is at an endpoint. It remains
to maximize on the three edges $\theta=4A$ for $A\in[0,1/16]$,
$\theta=8A$ for $A\in[0,1/32]$, and
$\theta=1/4$ for $A\in[1/32,1/16]$.
On each diagonal edge the risk has degree at most three in $A$, and on the
horizontal edge degree at most two. The certificate computes their exact
rational coefficients and verifies that every edge's second derivative
exceeds respectively $30$, $124$, and $3/100$ at both endpoints.
Since each such second derivative
is affine, it is positive throughout the edge. Thus each edge is convex
and the maximum over the entire triangle occurs at one of its three vertices.
This is a polynomial proof, not inference from a numerical grid.

The largest of the three exact vertex risks is less than $15443/10^6$;
its approximate value is $0.015442586588172843$. This proves the full-class
upper bound.

For the lower bound, use only the null $(0,0)$ and hard alternative
$(0,1/2)$, with probabilities $1-\pi,\pi$. Both belong to the full class.
For either fixed schedule $P\in\{O,N\}$ its optimal Bayes risk is
\[
 \mathcal B_P=\frac1{16}\pi(1-\pi)\sum_x
 \frac{P(0)_xP(1/32)_x}{(1-\pi)P(0)_x+\pi P(1/32)_x}.
\]
The preceding dominance argument for noninteractive cap-two designs
applies unchanged: the null has full finite-word support, so conditioning
on all seeds permits either one local policy with two flips or at most two
local policies with one flip each. Completing any early stops gives one
of these schedules and cannot increase the optimal Bayes risk. The
comparison permits arbitrary nonlinear estimates; the posterior mean used
in the Bayes formula is optimal among all of them. Therefore every
noninteractive procedure's worst-case MSE over the full class is at least
$\min\{\mathcal B_O,\mathcal B_N\}$.

The exact certificate verifies
\[
 \max_{\text{three vertices}} R
 <\frac{15443}{10^6}<\frac{15475}{10^6}
 <\min\{\mathcal B_O,\mathcal B_N\}.
\]
The smaller Bayes lower bound is approximately
$0.015475854685963823$. This proves the strict separation.
\end{proof}

All probabilities, coefficients, outputs and inequalities are recorded in
\linebreak\texttt{full\_class\_certificate.json}, generated by
\texttt{full\_class\_certificate.py} using rational arithmetic only.
The design prior was selected by a bounded exploratory rational bisection;
its optimality is not asserted or used. The independent proof of the full
triangle maximum is the essential new certification step.

This result shows that acquisition feedback can improve minimax performance
even when the adversary ranges over \emph{all} globally rare sign
allocations for a known magnitude law. It does not show that the improvement
grows with rarity or gap-scale count. The comparison for a uniform expected
budget remains separate, since a noninteractive procedure allowed occasional
work above two is not covered by the cap-two dominance argument.
